\begin{document}

\title{Heat Diffusion Solver}

\section{Types of Equations}
\subsection{Explicit}

\section{Parallel Architecture}
\subsection{Shared Memory / Intra-process / OpenMP}
\paragraph
    Updating each location in the domain discretization is performed in parallel.
    In the case of calculating average delta, a reduce is performed on the total.

\subsection{Multiple Memory / Inter-Process / OpenMPI}
\paragraph
    At the beginning of the simulation, a scatter is performed assigning partitions of the discrete domain to each process.
    Neighborhood communication groups are created for each process comprised of it and processes which own partitions adjacent to its partition.
    Note: We could describe this process with set builder notation.
\paragraph
    Each iteration step, a process broadcasts the outermost locations in its partition to its neighborhood.
    Likewise, the process async receives relevant locations from its neighbors.
    When a process uses locations it is to receive from its neighbors, it waits for the send to complete.
    It should perform these operations last.
\paragraph
    In the case of FDM, the domain is decomposed as a grid structure $G$ of shape $n_1\cross n_2 \cross ... \cross n_d$ where $d$ is the dimensionality.
    Each process is assigned cells in the partition
    \[
        G[0:n_1, 0:n_2, ...,r\lfloor\frac{n_d}{s}\rfloor:(r+1)\lfloor\frac{n_d}{s}\rfloor]
        \text{where } r = \text{rank of the process}
        \text{where } s = \text{world size (total processes)}
    \]
    In the implementation, the ownership of the last process will explicitly always end at $n_d$ to account for rounding errors, but that has been ommitted here for clarity.
    Note that this gaurantees cells are contiguous.
\paragraph
    Every X iteration step or when some global condition is met, processes syncronize their results to save them.
    (NOTE: This is subject to change I don't know what HDF5 is precisely but it may fix this.)

\end{document}
